\documentclass[../../main.tex]{subfiles}

\begin{document}

\section{Introduction to the density and the reduced density matrix}
\label{sec:density-and-reduced-density-matrix}

\subsection{The density operator}

Let $\{p_i,\ket{\psi_i}\}$ be an ensemble of pure states, describing a system that is in one
of the states $\ket{\psi_i}$ with probability $p_i$, i.e.\ the system is not fully known. Define
the density operator as
\begin{equation}
\rho=\sum_i p_i\,\dyad{\psi_i}{\psi_i},\qquad p_i\in\mathbb R\ \, \&\ \sum_i p_i=1,
\end{equation}
which is a positive semi-definite operator of trace one (see for example Nielsen \& Chuang).

\subsection{Observables on a composite system}

Take a system $A$ with some observable $M$, and a system $B$, giving the composite system
$AB$ (see Nielsen \& Chuang Box 2.6). Assume we can measure $M$ on system $A$; then a corresponding measurement
on $AB$ is $\tilde M=M\otimes\mathds{1}_B$.

\begin{proof}
Take $AB$ to be prepared in $\ket{m}\ket{\psi}$ with $M\ket{m}=m\ket{m}$. Applying the
measurement to $\ket{m}\ket{\psi}$ should yield $m$ with probability one, i.e.\ the
corresponding projector for $\tilde M$ is $P_m\otimes\mathds{1}_B$. By the spectral theorem,
$\tilde M=\sum_m m\,P_m\otimes\mathds{1}_B=M\otimes\mathds{1}_B$ for a general state of $AB$.
\end{proof}

\subsection{The reduced density matrix as a partial trace}

Measurement averages must be the same whether computed via the full $\rho^{AB}$ or via some
$\rho^{A}$ associated to system $A$, i.e.
\begin{equation}
\Trace\bigl(M\rho^A\bigr)\overset{!}{=}\Trace\bigl(\tilde M\rho^{AB}\bigr)
=\Trace\bigl((M\otimes\mathds{1}_B)\rho^{AB}\bigr).
\end{equation}
Take $\rho^A=\Trace_B(\rho^{AB})$ where
\begin{equation}
\Trace_B\bigl(\dyad{a_1}{a_2}\otimes\dyad{b_1}{b_2}\bigr)
\equiv\dyad{a_1}{a_2}\,\Trace\bigl(\dyad{b_1}{b_2}\bigr)
\qquad\&\qquad\Trace_B(\cdot)\ \text{linear.}
\end{equation}
Write $\rho^{AB}$ by using the eigenbasis of $M$ for system $A$ with $M\ket{i}=m_i\ket{i}$ and
$\braket{i}{j}=\delta_{ij}$:
\begin{align}
\rho^{AB}&=\sum_\alpha p^\alpha\dyad{\psi}{\psi}^\alpha
=\sum_\alpha p^\alpha\sum_{ijk\ell}c_{ij}^\alpha(c_{k\ell}^\alpha)^*\,
\ket{i}_A\ket{j}_B\bra{\ell}_B\bra{k}_A,\\
\rho^{A}&=\sum_\alpha\sum_{ijk\ell}p^\alpha c_{ij}^\alpha(c_{k\ell}^\alpha)^*\,
\dyad{i}{k}_A\underbrace{\Trace\bigl(\dyad{j}{\ell}_B\bigr)}_{=\delta_{j\ell}}
=\sum_{\alpha,ijk}p^\alpha c_{ij}^\alpha(c_{kj}^\alpha)^*\,\dyad{i}{k}_A.
\end{align}
Therefore
\begin{align}
\Trace\bigl(M\rho^A\bigr)
&=\sum_{\alpha,ijk}p^\alpha c_{ij}^\alpha(c_{kj}^\alpha)^*\,
\underbrace{\Trace\bigl(m_i\dyad{i}{k}_A\bigr)}_{=m_i\delta_{ik}}
=\sum_{\alpha,ij}p^\alpha m_i\,c_{ij}^\alpha(c_{ij}^\alpha)^*,
\end{align}
and similarly $\Trace\bigl((M\otimes\mathds{1}_B)\rho^{AB}\bigr)
=\sum_{\alpha,ij}p^\alpha m_i\,c_{ij}^\alpha(c_{ij}^\alpha)^*$, which proves the claim. One can
also easily show that this is the only map that conserves consistency (see Nielsen \& Chuang Box 2.6).

\subsection{Reduced density matrix of a pure state}

Now examine the reduced density matrix for a pure state,
\begin{align}
\rho_{AB}&=\dyad{\psi}{\psi}
=\sum_{ijk\ell}c_{ij}(c_{k\ell})^*\,\ket{i}_A\ket{j}_B\bra{\ell}_B\bra{k}_A,\\
\rho_A&=\Trace_B(\rho_{AB})=\sum_{ijk}c_{ij}(c_{kj})^*\,\dyad{i}{k}_A, \label{eq:reduced_density_matrix_as_operator}
\end{align}
such that the matrix representation is
\begin{equation}
\boxed{\,\rho_A=\psi\psi^\dagger\,}, \label{eq:density_matrix_as_wavefunction_coefficients}
\end{equation}
using that the $c_{ij}$ parametrize the matrix $\psi\in\mathbb C^{d_A\times d_B}$, i.e.\
$\ket{\psi}=\sum_{ij}c_{ij}\ket{i}_A\ket{j}_B$ with $\sum_{ij} \vert c_{ij} \vert^2 = 1$.


Furthermore, \cref{eq:reduced_density_matrix_as_operator} results in
\begin{align}
      \Trace \rho_A &= \Trace \left[ \sum_{ijk} c_{ij} (c_{kj})^* \dyad{i}{k}_A \right] \notag \\
      &= \sum_{ijkl} c_{ij} (c_{kj})^* \delta_{il} \delta_{kl} \notag \\
      &= \sum_{jl} c_{lj} (c_{lj})^* = 1 , \notag
\end{align}
which means that that the eigenvalues of the reduced density matrix sum to one. This can be seen using that any hermitian matrix can be diagonalized and the cylic property of the trace
\begin{equation}
      \Trace \rho_A = \Trace [O D O^\dagger] = \Trace D = \sum_i \lambda_i = 1. \label{eq:eigenvalues_reduced_density_matrix_sum_to_one}
\end{equation}


\subsection{Connection to the SVD}

Now use the SVD on $\psi$:
\begin{equation}
\psi=USV^\dagger\qquad\text{with }U\in\mathbb C^{d_A\times d_A},\ V\in\mathbb C^{d_B\times d_B}\ \text{unitary.}
\end{equation}
Then, using $V^\dagger V=\mathds 1$,
\begin{equation}
\rho_A=USV^\dagger VS^\dagger U^\dagger
=USS^\dagger U^\dagger,
\qquad
SS^\dagger=\diag(s_1^2,\dots,s_r^2,0,\dots,0),
\end{equation}
and write $S^2 = \diag(s_1^2,\dots,s_r^2)$.
Furthermore, since $\rho_A$ is hermitian we can diagonalize it with
$\rho_A=ODO^\dagger=US^2 U^\dagger$ and $D=\diag(\lambda_1,\dots)$, so that
\begin{equation}
\boxed{\,\lambda_i=s_i^2\,},
\end{equation}
i.e.\ the eigenvalues are given by the squared Schmidt values and hence non-negative. Thus, we can conclude together with \cref{eq:eigenvalues_reduced_density_matrix_sum_to_one} that 
\begin{equation}
      \boxed{\, 0 \leq \lambda_i \leq 1 \,} . \label{eq:range_of_eigenvalues_of_reduced_density_matrix}
\end{equation}




\subsection{References}
\begin{itemize}
\item M.\ Nielsen \& I.\ Chuang, \emph{Quantum Computation and Quantum Information}
      (density operator, partial trace; Box 2.6).
\item U.\ Schollw\"ock, \emph{The density-matrix renormalization group in the age of matrix
      product states}, beginning of Chapter 4 (MPS).
\end{itemize}

\end{document}
